%-----------------------------------------------------------------------
% Polytechnic University of Bari
% Year's Work Report or Graduation Thesis Model in LaTEX
% Prof. Eng. Francesco Adamo, PhD
% October 2021
%
% Based on the work of Loïc Barrault, Le Mans Universitè - September 2021
% and on the work of Prof. Roger K. Moore, University of Sheffield - July 2018
%------------------------------------------------------------------------

\documentclass[12pt,twoside,italian]{book} % se necessario cambiare il vocabolario qui
\usepackage[bindingoffset=1cm, margin=2cm]{geometry} % offset di 1 cm necessario per garantire una corretta rilegatura del volume
\usepackage[toc,page]{appendix}
\usepackage{graphicx}
\usepackage{lipsum}
\usepackage[hidelinks]{hyperref} % opzione hidelinks per sopprimere i box rossi intorno ai link ipertestuali negli indici
\usepackage[utf8]{inputenc}
\usepackage[T1]{fontenc}
\usepackage[italian]{babel} % se necessario cambiare il vocabolario qui 
\usepackage{caption}
\usepackage{nextpage} % aggiunge il comando \cleartooddpage[1]
\usepackage{pdfpages}
\usepackage{setspace}
\usepackage{ragged2e}
\usepackage{xcolor}
\usepackage{array}
\usepackage{gensymb}
\usepackage{textcomp} 
\usepackage[normalem]{ulem}
\usepackage{titlesec}
\usepackage[numbers]{natbib} % per ottenere la bibliografia e i collegamenti con stile numerato
\usepackage{listings}
\usepackage{float}

% Definisce lo stile di numerazione per i titoli dei paragrafi di quarto livello
\titleclass{\subsubsubsection}{straight}[\subsection]

\newcounter{subsubsubsection}[subsubsection]
\renewcommand\thesubsubsubsection{\thesubsubsection.\arabic{subsubsubsection}}

\titleformat{\subsubsubsection}
  {\normalfont\normalsize\bfseries}{\thesubsubsubsection}{1em}{}
\titlespacing*{\subsubsubsection}
{0pt}{3.25ex plus 1ex minus .2ex}{1.5ex plus .2ex}
\makeatletter
\def\toclevel@subsubsubsection{4}
\def\l@subsubsubsection{\@dottedtocline{4}{7em}{4em}}
\makeatother
\setcounter{secnumdepth}{4}
\setcounter{tocdepth}{4}

% Modifica titoli predefiniti per le appendici
\renewcommand{\appendixtocname}{Appendici}%
\renewcommand{\appendixpagename}{Appendici}%

\begin{document}

\captionsetup[figure]{margin=1.5cm,font=small,labelfont={bf},name={Figura},labelsep=colon,textfont={it}}
\captionsetup[table]{margin=1.5cm,font=small,labelfont={bf},name={Tabella},labelsep=colon,textfont={it}}
\SetLipsumDefault{1}

% -------------------------------------------------------------------
% Frontespizio
% -------------------------------------------------------------------
\addcontentsline{toc}{chapter}{frontespizio}
\input{sections/frontespizio.tex}

% -------------------------------------------------------------------
% Modulo della liberatoria alla consultazione della tesi
% -------------------------------------------------------------------
\cleartooddpage
\includepdf[pages=1,pagecommand={\phantomsection\addcontentsline{toc}{chapter}{liberatoria}}]{sections/modulo_liberatoria.pdf} % eventualmente sostituire il file con la scansione del modulo compilato

% -------------------------------------------------------------------
% Dedica
% -------------------------------------------------------------------
\input{sections/dedica.tex}

% -------------------------------------------------------------------
% Ringraziamenti
% -------------------------------------------------------------------
\input{sections/ringraziamenti.tex}


% Ridefinisce i nomi delle sezioni dell'indice delle figure e dell'indice delle tabelle
\renewcommand{\listfigurename}{Indice delle figure}
\renewcommand{\listtablename}{Indice delle tabelle}

% -------------------------------------------------------------------
% Sommario, indice delle figure e indice delle tabelle
% -------------------------------------------------------------------
\tableofcontents
\addcontentsline{toc}{chapter}{Indice}

\listoffigures
\addcontentsline{toc}{chapter}{Indice delle figure}

\listoftables
\addcontentsline{toc}{chapter}{Indice delle tabelle}


% -------------------------------------------------------------------
% Introduzione generale (o abstract)
% -------------------------------------------------------------------
\input{sections/introduzione.tex}

% -------------------------------------------------------------------
% Contenuto principale
% -------------------------------------------------------------------
\mainmatter
\input{sections/capitolo_1/capitolo_1.tex}
\input{sections/capitolo_2/capitolo_2.tex}
\input{sections/conclusioni.tex}
\input{sections/bibliografia.tex}

% -------------------------------------------------------------------
% Appendici
% -------------------------------------------------------------------

\begin{appendices}
\input{sections/appendice_A/appendice_A.tex}
\end{appendices}

\end{document}
